% Part:sets-functions-relations
% Chapter: size-of-sets
% Section: zig-zag

\documentclass[../../../include/open-logic-section]{subfiles}

\begin{document}

\olfileid[pt]{sfr}{siz}{zigzag}

\olsection{Método do Zigue-Zague de Cantor}

\begin{explain}
Já consideramos algumas enumerações ``fáceis''. Vejamos algo um pouco mais
difícil. Considere-se o conjunto dos pares de números
naturais\oliflabeldef{sfr:set:pai:sec}{, que definimos na
\olref[set][pai]{sec} assim:}{definido por:}
\[
\Nat \times \Nat = \Setabs{\tuple{n,m}}{n,m \in \Nat}
\]
Podemos organizar estes pares ordenados numa \emph{matriz}, assim:
\[
\begin{array}{ c | c | c | c | c | c}
& \mathbf 0 & \mathbf 1 & \mathbf 2 & \mathbf 3 & \dots \\
\hline
\mathbf 0 & \tuple{0,0} & \tuple{0,1} & \tuple{0,2} & \tuple{0,3} & \dots \\
\hline
\mathbf 1 & \tuple{1,0} & \tuple{1,1} & \tuple{1,2} & \tuple{1,3} & \dots \\
\hline
\mathbf 2 & \tuple{2,0} & \tuple{2,1} & \tuple{2,2} & \tuple{2,3} & \dots \\
\hline
\mathbf 3 & \tuple{3,0} & \tuple{3,1} & \tuple{3,2} & \tuple{3,3} & \dots \\
\hline
\vdots & \vdots & \vdots & \vdots & \vdots & \ddots\\
\end{array}
\]
Claramente, cada par ordenado em $\Nat \times \Nat$ aparece exatamente
uma vez na matriz. Em particular, $\tuple{n,m}$ aparece na $n$-ésima linha e na
$m$-ésima coluna. Mas como organizar os elementos de tal matriz numa lista
``unidimensional''? O padrão na matriz abaixo demonstra uma forma de o fazer (embora
haja, claro, muitas outras):
\[
\begin{array}{ c | c | c | c | c | c | c}
& \mathbf 0 & \mathbf 1 & \mathbf 2 & \mathbf 3 & \mathbf 4 &\dots \\
\hline
\mathbf 0 & 0  & 1& 3 & 6& 10 &\ldots \\
\hline
\mathbf 1 &2 & 4& 7 & 11 & \dots &\ldots \\
\hline
\mathbf 2 & 5 & 8 & 12 & \ldots & \dots&\ldots \\
\hline
\mathbf 3 & 9 & 13 & \ldots & \ldots & \dots & \ldots \\
\hline
\mathbf 4 & 14 & \ldots & \ldots & \ldots & \dots & \ldots \\
\hline
\vdots & \vdots & \vdots & \vdots & \vdots&\ldots & \ddots\\
\end{array}
\]\noindent
Este padrão chama-se \emph{método do zigue-zague de Cantor}. Enumera
$\Nat \times \Nat$ do seguinte modo:
\[
\tuple{0,0}, \tuple{0,1}, \tuple{1,0}, \tuple{0,2}, \tuple{1,1},
\tuple{2,0}, \tuple{0,3}, \tuple{1,2}, \tuple{2,1}, \tuple{3,0}, \dots
\]
Isto estabelece o seguinte:
\end{explain}

\begin{prop}\ollabel{natsquaredenumerable}
$\Nat \times \Nat$ é !!{enumerable}.
\end{prop}

\begin{proof}
Seja $f \colon \Nat \to \Nat\times\Nat$ a função que a cada $k \in \Nat$ associa o
tuplo $\tuple{n,m} \in \Nat \times \Nat$ tal que $k$ é o valor na $n$-ésima
linha e $m$-ésima coluna da matriz do zigue-zague de Cantor.
\end{proof}

\begin{explain}
Esta técnica generaliza-se bem. Por exemplo, podemos usá-la para enumerar o
conjunto dos triplos ordenados de números naturais, isto é:
\[
\Nat \times \Nat \times \Nat = \Setabs{\tuple{n,m,k}}{n,m,k \in \Nat}
\]
Pensamos em $\Nat \times \Nat \times \Nat$ como o produto cartesiano de $\Nat
\times \Nat$ com $\Nat$, isto é,
\[
\Nat^3 = (\Nat \times \Nat) \times \Nat =
\Setabs{\tuple{\tuple{n,m},k}}{n, m, k
  \in \Nat }
\]
e assim podemos enumerar $\Nat^3$ com uma matriz rotulando um eixo com a
enumeração de~$\Nat$ e o outro com a enumeração de~$\Nat^2$:
\[
\begin{array}{ c | c | c | c | c | c}
& \mathbf 0 & \mathbf 1 & \mathbf 2 & \mathbf 3 & \dots \\
\hline
\mathbf{\tuple{0,0}} & \tuple{0,0,0} & \tuple{0,0,1} & \tuple{0,0,2} & \tuple{0,0,3} & \dots \\
\hline
\mathbf{\tuple{0,1}} & \tuple{0,1,0} & \tuple{0,1,1} & \tuple{0,1,2} & \tuple{0,1,3} & \dots \\
\hline
\mathbf{\tuple{1,0}} & \tuple{1,0,0} & \tuple{1,0,1} & \tuple{1,0,2} & \tuple{1,0,3} & \dots \\
\hline
\mathbf{\tuple{0,2}} & \tuple{0,2,0} & \tuple{0,2,1} & \tuple{0,2,2} & \tuple{0,2,3} & \dots\\
\hline
\vdots & \vdots & \vdots & \vdots & \vdots & \ddots \\
\end{array}
\]
Assim, com um método como o zigue-zague de Cantor, obtemos analogamente uma
enumeração de~$\Nat^3$. E podemos continuar, obtendo enumerações de $\Nat^n$ para
qualquer número natural~$n$. Logo:
\end{explain}

\begin{prop}
$\Nat^n$ é !!{enumerable}, para todo $n \in \Nat$.
\end{prop}

\begin{prob}\label{sfr:siz:zigzag:prob:posint-n}
Mostre que $(\PosInt)^n$ é !!{enumerable}, para todo $n \in \Nat$.
\end{prob}

\begin{prob}\label{sfr:siz:zigzag:prob:posint-star}
  Mostre que $(\PosInt)^*$ é !!{enumerable}. Pode assumir \cref{sfr:siz:zigzag:prob:posint-n}.
\end{prob}

\end{document}
